% paper/sections/03b_cls_transport.tex
\subsection{保存形移流方程式}
\label{sec:cls_advection}

§~\ref{sec:why_cls} で示した質量保存問題を解決するため，CLS 法（\cite{OlssonKreiss2005}；
古典 Level Set 文脈の総覧は~\cite{OsherFedkiw2003}）では
界面指標 $\psi(\bx, t)$（平滑化ヘヴィサイド関数）の移流方程式を保存形で記述する：
\begin{equation}
  \pd{\psi}{t} + \bnabla\cdot(\psi\,\bu) = 0
  \label{eq:cls_advection}
\end{equation}
本稿の符号規約では $\psi\approx1$ が液相，$\psi\approx0$ が気相であるため，
閉じた計算領域で $\int \psi \, dV$（液相体積に対応）が保存されれば，
液相体積が直接保存される．

\begin{tcolorbox}[resultbox, title={保存形の体積保全性：発散定理の直接適用と離散精度}]
保存形移流方程式~\eqref{eq:cls_advection} を領域 $\Omega$ で積分し，
Gauss の発散定理を $\bnabla\cdot(\psi\bu)$ に直接適用すると：
\[
  \frac{d}{dt}\int_\Omega\psi\,dV
  = -\int_\Omega\bnabla\cdot(\psi\bu)\,dV
  = -\oint_{\partial\Omega}\psi\bu\cdot\hat{\bm n}_{\partial\Omega}\,dS
\]
ここで $\hat{\bm n}_{\partial\Omega}$ は計算領域境界 $\partial\Omega$ の外向き法線（§~\ref{sec:notation} で定義した界面法線 $\hat{\bm n}_{lg}$ とは別の幾何対象）である．
これは連続レベルでは $\bnabla\cdot\bu$ の値に無関係に成立する．
連続レベルでは，壁面境界（$\bu\cdot\hat{\bm n}_{\partial\Omega}=0$）または周期境界では右辺がゼロとなり，
$\int\psi\,dV = \text{const}$ が速度場の発散条件に依存せず保証される．
離散計算では，この性質を保つ精度はフラックス差分，境界処理，値域制限，再初期化補正に依存する．

さらに，静止 no-slip 壁に接する界面では，体積保存だけでなく壁面上の相トポロジも制約される．
壁面を接線座標 $s$ で表し $\psi_w(s,t)=\psi(\bx_{\partial\Omega}(s),t)$ とおくと，
壁面上では $\bu=0$ であるため，明示的な slip/contact-line/contact-angle モデルを導入しない限り，
\begin{equation}
  \operatorname{sign}\!\left(\psi_w(s,t)-\frac12\right)
  =
  \operatorname{sign}\!\left(\psi_w(s,0)-\frac12\right),
  \qquad
  C(t)=C(0),\quad
  C(t)=\{s:\psi_w(s,t)=1/2\}
  \label{eq:wall_phase_topology_invariant}
\end{equation}
である．
これは平滑化された $\psi_w$ の数値を全時刻で固定する条件ではない：
界面が壁に触れない範囲で壁近傍へ近づけば，同じ相側に留まる $\psi_w$ の値は変化しうる．
固定されるのは壁面上の相区間と既存の接触点であり，再初期化・格子再構成・質量補正などの補助操作は
新しい壁面接触点を生成してはならない．
周期境界には物理的な壁面接触集合 $C(t)$ が存在しないため，この制約は適用しない．

これが非保存形（式~\eqref{eq:ls_noncons}）との本質的な差異である：
非保存形では $\bu\cdot\bnabla\phi \equiv \bnabla\cdot(\phi\bu) - \phi\bnabla\cdot\bu$
の等価性に $\bnabla\cdot\bu=0$ が必要であり，
離散発散 $[\bnabla\cdot\bu]_h\neq 0$ が体積ドリフトに直結する（§~\ref{sec:why_cls_massloss}）．

\textbf{離散化における精度：}
\begin{itemize}
  \item 本稿の標準経路では CLS 移流に FCCD 保存形面フラックス
        （§~\ref{sec:advection_motivation}）を使用する．
        共有面フラックスの差分として $\bnabla\cdot(\psi\bu)$ を組むため，
        内部面の寄与は対で相殺される．残る体積残差は境界フラックス，
        値域制限，格子転写，再初期化後の質量閉鎖に分離して監視する
        （§~\ref{sec:cls_stages}）．
  \item 圧力 Poisson 求解前の中間速度 $\bu^*$（本稿のアルゴリズム構成における予測段出力）を CLS 移流に用いると，
        壁面法線速度が厳密にゼロでない場合に境界フラックスへの誤差が生じうる．
        標準経路では，前ステップの投影で得た射影後の面速度 $\bu_f^n$
        を式~\eqref{eq:projection_native_psi_transport} の CLS フラックスへ直接渡す．
  \item 第~\ref{sec:verification}章の検証で，静止・密度比・純移流の各条件における
        $\int\psi\,dV$ 残差を定量確認する．
\end{itemize}

\begin{tcolorbox}[defbox, title={共通フラックス台帳：界面だけを保存しても足りない場合}]
保存形 CLS の面フラックスを
\[
  F_q^r=(P_f q^r)u_f^r
\]
と記録する．ここで $r$ は TVD-RK3 の段階である．
密度比が小さい受動界面では，この台帳は主に液相体積の監視に使えばよい．
一方，水--空気級の上昇気泡のように密度比が大きく，
界面移流そのものが運動量輸送を支配する場合は，同じ台帳から
\[
  F_m^r=\rho_g F_V^r+(\rho_l-\rho_g)F_q^r,\qquad
  F_{\bm{p}}^r=F_m^r\bu_{\mathrm{up}}^r
\]
を作る．$F_V^r$ は同じ面速度から得る幾何体積フラックスであり，
$\bu_{\mathrm{up}}^r$ は運動量に用いる上流速度である．
このとき，更新は
\[
  q^{r+1}=q^r-\Delta t\,D_fF_q^r,\quad
  m^{r+1}=m^r-\Delta t\,D_fF_m^r,\quad
  \bm{p}^{r+1}=\bm{p}^r-\Delta t\,D_fF_{\bm{p}}^r
\]
を同じ離散発散 $D_f$ で満たさなければならない．
圧力射影で発散ゼロと判定した面速度を，別の発散演算子で
界面輸送へ渡すと，射影の意味で非圧縮でも輸送の意味で圧縮的な速度場になりうる．
したがって台帳には，段階ごとの $q^r$，面フラックス，使用した $D_f$，
値域制限の有無を残す．
端点の体積だけが合っていることは，共通フラックスの証明にはならない．
\end{tcolorbox}

\medskip\noindent\textbf{CLS が標準 LS より優れる理由：保存責務の分離}

コロケート格子では MAC スタガード格子と異なり
「速度場そのものの厳密な離散発散ゼロ性」は構造的には保証されない．
本稿の CLS 経路の優位性は，その不足を $\phi$ 非保存形移流へ隠すのではなく，
$\psi$ の共有面フラックス，境界処理，値域制限，質量閉鎖へ分解して監視できる点にある：

\begin{center}
\renewcommand{\arraystretch}{\PaperTableArrayStretch}
\begin{tabular}{p{3.5cm}p{4.5cm}>{\raggedright\arraybackslash}p{5cm}}
\hline
方法 & 保存性の扱い & 主な誤差源・監視対象 \\
\hline
標準 LS（非保存形） & 非圧縮条件との離散不整合が体積ドリフトへ入る &
  $\phi[\bnabla\cdot\bu]_h$ の非キャンセル項 \\
CLS（保存形 + FCCD 共有面フラックス） & 内部面フラックスは対消滅し，残差を境界・値域制限・再初期化・質量閉鎖へ分離 &
  境界フラックス，値域制限，格子転写，段階 B/F 質量閉鎖残差 \\
MAC スタガード格子 & 保存形フラックスと離散発散ゼロ速度を組み合わせやすい &
  スカラー移流の打ち切り誤差，値域制限，境界処理 \\
\hline
\end{tabular}
\end{center}

標準 LS は非保存形のため，Projection 法の PPE 空間離散化に起因する離散発散残差
$[\bnabla\cdot\bu]_h = \Ord{h^2}$（2次精度 PPE の場合）が
$\phi\,[\bnabla\cdot\bu]_h$ として体積誤差に直結する（§~\ref{sec:why_cls_massloss}）．
CLS（保存形 + FCCD）では保存形の発散定理に対応する共有面フラックス差分により，
この非キャンセル項を $\psi$ 輸送から排除する．
したがって体積保存の議論は，フラックス発散そのものではなく，境界，値域制限，
再初期化，および段階 B/F の質量閉鎖残差として評価される．
\end{tcolorbox}

\subsection{再初期化方程式（等方拡散形）}
\label{sec:reinit}

移流によって $\psi$ のプロファイルが崩れるため，界面幅と単調性を保つ幾何修復が必要になる．
古典的 CLS では，目標平衡プロファイルを以下の圧縮・拡散方程式の定常解として特徴づける：
\begin{equation}
 \pd{\psi}{\tau} + \bnabla\cdot\bigl(\psi(1-\psi)\hat{\bm n}_c\bigr) = \bnabla\cdot\bigl(\varepsilon\,\bnabla\psi\bigr)
 \label{eq:cls_reinit_iso}
\end{equation}

ここで圧縮法線は $\hat{\bm n}_c=\bnabla\psi/|\bnabla\psi|$ とし，
本稿の幾何法線 $\hat{\bm n}_{lg}=\bnabla\phi/|\bnabla\phi|$ とは
$\hat{\bm n}_c=-\hat{\bm n}_{lg}$ の関係にある．
ただし CLS の基本変数は $\psi$ であり，$\phi$ は $\psi$ からロジット逆変換
$\phi = \varepsilon\ln\!\bigl((1-\psi)/\psi\bigr)$ で再構成する
（§~\ref{sec:newton_inversion}）．
飽和帯（$\psi < \psi_{\min}$ または $\psi > 1-\psi_{\min}$，$\psi_{\min} = 0.01$）では
$\psi$ 勾配とロジット引数が数値退化するため，曲率は界面帯だけで評価し，
$\phi$ が必要な HFE・壁面幾何・距離閉鎖では品質監視後に
Ridge--Eikonal 距離再構成（§~\ref{sec:ridge_eikonal}）を用いる．
これは低次の代替スキームではなく，標準 CLS 手順の幾何修復段である．

式~\eqref{eq:cls_reinit_iso} は「定常解として $\psi = H_\varepsilon(-\phi)$ を持つように逆算して構成した」という
設計原理に基づく．圧縮項 $\bnabla\cdot(\psi(1-\psi)\hat{\bm n}_c)$ と拡散項 $\bnabla\cdot(\varepsilon\bnabla\psi)$ が
ちょうど釣り合う不動点が $H_\varepsilon$ であることの代数的証明は付録~\ref{app:cls_fixed_point} を参照のこと．

\textbf{界面位置（$\psi=0.5$ 等値面）の準不動性．}
圧縮項と拡散項を展開すると，$\psi=0.5$ では
法線方向の勾配成分に現れる係数 $(1-2\psi)$ と
定常プロファイル $H_\varepsilon$ の変曲点条件 $H_\varepsilon''(0)=0$ が効く．
一方で $\bnabla\cdot\hat{\bm n}_c$ に比例する曲率項は残るため，
界面位置の不動性は「圧縮項単独がゼロ」ではなく，
圧縮項と拡散項が定常プロファイル上で釣り合うことにより成立する．
その結果，定常状態では $\partial\psi/\partial\tau\big|_{\psi=0.5} = 0$ が成立する
（代数的証明：付録~\ref{app:cls_fixed_point}）．
したがって本稿では，圧縮--拡散式を「どのプロファイルへ戻すべきか」を定める
連続レベルの平衡条件として用い，離散標準経路では段階 D/F の
距離再構成と質量閉鎖で零点集合移動量を監視する（詳細：付録~\ref{app:cls_fixed_point}）．

% ─────────────────────────────────────────────────────────────────────────────
\subsubsection{$\varepsilon$ のスケーリングと界面幅}
% ─────────────────────────────────────────────────────────────────────────────

右辺の拡散項係数 $\varepsilon$ は，数値的な界面幅を決定するスケールパラメータである．
数値計算上はグリッド解像度に応じた有限の界面幅が必要であり，以下のスケーリングを用いる：
\begin{equation}
  \varepsilon = C_\varepsilon \, \Delta x_{\min} \quad (C_\varepsilon \approx 1.0 \sim 2.0)
\end{equation}

\paragraph{非一様格子における $\Delta x_{\min}$ の定義}
非一様格子を使う場合，
$\Delta x_{\min}$ は界面が現在存在する領域（$|\phi| \lesssim 3\varepsilon$ 付近）での
局所最小格子幅として設定する：
\[
  \Delta x_{\min}
  \;=\; \min_{(i,j):\,|\phi_{i,j}| \le 3\varepsilon}\!\bigl(\Delta x_{i,j},\,\Delta y_{i,j}\bigr)
\]
$\varepsilon$ は再初期化方程式の拡散係数として機能し，
その拡散効果が有意な範囲は $|\phi| \lesssim 3\varepsilon$ に限定される
（$\psi_{\mathrm{eq}}(\phi) = H_\varepsilon(-\phi)$ が $|\phi| > 3\varepsilon$ で飽和するため
$\partial\psi/\partial\tau \approx 0$）．
遠方の粗格子点では $\psi(1-\psi) \approx 0$ となり再初期化の影響を受けない．
なお，界面が移動するシミュレーションでは
$\Delta x_{\min}$ を毎タイムステップで界面位置に基づいて更新する必要がある．

CLS 変数 $\psi$ の定常プロファイルは，$\phi$ を用いた滑らかなヘヴィサイド関数として定義される：
\begin{equation}
  \psi_{\text{eq}}(\phi) = H_\varepsilon(-\phi) = \frac{1}{1 + e^{\phi/\varepsilon}}
\end{equation}
この関数は $\phi \in [-3\varepsilon, +3\varepsilon]$ の範囲で $\psi$ が 1 から 0 へと滑らかに遷移する．
本稿では中心遷移幅を約 $2\varepsilon$，95\% 遷移帯を約 $6\varepsilon$ と呼び分ける．

% ─────────────────────────────────────────────────────────────────────────────
\subsubsection{平滑化デルタ関数 \texorpdfstring{$\delta_\varepsilon$}{δ\_ε} と界面近似}
\label{sec:smoothed_delta}
% ─────────────────────────────────────────────────────────────────────────────

$H_\varepsilon(-\phi)$ の微分は閉形式で得られ，界面近傍にのみ集中するピーク関数を与える：
\begin{equation}
  \delta_\varepsilon(\phi) \equiv H_\varepsilon'(-\phi)
  = \frac{1}{\varepsilon}H_\varepsilon(-\phi)\bigl(1-H_\varepsilon(-\phi)\bigr)
  = \frac{1}{\varepsilon}\psi(1-\psi)
  \label{eq:delta_eps_def}
\end{equation}
$d\psi/d\phi=-\delta_\varepsilon$ である．
$\delta_\varepsilon$ は偶関数（$\delta_\varepsilon(\phi) = \delta_\varepsilon(-\phi)$）であり，
$\phi=0$（界面）で最大値 $1/(4\varepsilon)$ をとり，$|\phi| \gg \varepsilon$ で急速に減衰する．

\begin{figure}[htbp]
\centering
\begin{tikzpicture}[scale=1.0]
  %% 左：ヘヴィサイド関数
  \begin{scope}[xshift=-45mm]
    \node[font=\small\bfseries, text=navy] at (0, 2.8) {（a）ヘヴィサイド関数 $H_\varepsilon(-\phi)$};
    \draw[->] (-2.2,0) -- (2.2,0) node[right, font=\small]{$\phi$};
    \draw[->] (0,-0.3) -- (0, 2.5) node[above, font=\small]{$H_\varepsilon$};
    \draw[very thick, blue2, domain=-2:2, samples=80]
      plot (\x, {1.8/(1+exp(3*\x))});
    \node[font=\scriptsize, blue2] at (-1.4, 1.65) {$H_\varepsilon(-\phi)$};
    \draw[dashed, gray] (-2.2, 1.8) -- (2.2, 1.8) node[right, font=\scriptsize, gray]{$1$};
    \draw[dashed, gray] (-2.2, 0)   -- (2.2, 0)   node[right, font=\scriptsize, gray]{$0$};
    \node[font=\scriptsize, red2] at (0.15, -0.25) {$\phi{=}0$};
    \node[font=\scriptsize] at (-1.3, 0.9) {液相};
    \node[font=\scriptsize] at (1.3,  0.9) {気相};
  \end{scope}
  %% 右：デルタ関数（微分）
  \begin{scope}[xshift=45mm]
    \node[font=\small\bfseries, text=navy] at (0, 2.8) {（b）微分 $\delta_\varepsilon(\phi) = H_\varepsilon'(-\phi)$};
    \draw[->] (-2.2,0) -- (2.2,0) node[right, font=\small]{$\phi$};
    \draw[->] (0,-0.3) -- (0, 2.5) node[above, font=\small]{$\delta_\varepsilon$};
    \draw[very thick, red2, domain=-2:2, samples=80]
      plot (\x, {2.0*exp(-3*\x*\x)});
    \fill[red2!15, domain=-2:2, samples=80]
      plot (\x, {2.0*exp(-3*\x*\x)}) -- (2, 0) -- (-2, 0) -- cycle;
    \node[font=\scriptsize, red2] at (0.6, 1.8) {$\delta_\varepsilon(\phi)$};
    \node[font=\scriptsize, red2] at (0.15, -0.25) {$\phi{=}0$};
    \draw[<->, thick, darkgray]
      (-0.6, -0.15) -- (0.6, -0.15) node[midway, below, font=\scriptsize]{中心幅 $\approx 2\varepsilon$};
    \node[font=\scriptsize, darkgray] at (1.6, 0.3) {界面外 $\approx 0$};
    \node[font=\scriptsize] at (0, 1.0) {$\int\delta_\varepsilon\,d\phi = 1$};
  \end{scope}
  %% 微分矢印
  \draw[-Stealth, very thick, navy] (-0.8, 1.4) -- (0.8, 1.4)
        node[midway, above, font=\small, navy]{$-\dfrac{d}{d\phi}$};
\end{tikzpicture}
\caption{（a）CLS 変数 \texorpdfstring{$\psi = H_\varepsilon(-\phi)$}{psi = Heps(-phi)}（滑らかヘヴィサイド関数）}
\label{fig:heaviside_delta}
\PaperCaptionNote{界面 $\phi=0$ を中心に液相側の 1 から気相側の 0 へ連続的に変化する． （b）$\delta_\varepsilon(\phi) = H_\varepsilon'(-\phi)$ は界面近傍にのみ集中した ピーク（ディラックデルタ関数の滑らか版）を形成し，$d\psi/d\phi=-\delta_\varepsilon$ を満たす．}
\end{figure}

$\delta_\varepsilon$ は CLS 変数 $\psi$ の勾配と以下の関係にある：
\begin{equation}
  |\bnabla\psi| = |H_\varepsilon'(\phi)|\,|\bnabla\phi|
               = \delta_\varepsilon(\phi)\,|\bnabla\phi|
               \approx \delta_\varepsilon(\phi)
  \label{eq:grad_psi_delta}
\end{equation}
最後の近似は $|\bnabla\phi|\approx 1$（Eikonal 条件）による．
すなわち，\textbf{$|\bnabla\psi|$ が界面ディラックデルタ $\delta_s$ の滑らかな近似}となる．
この性質は，表面張力の体積力表現（CSF モデル，付録~\ref{app:csf_full}）や
HFE（§~\ref{sec:field_extension}；分相 PPE 解法の基盤技術）における界面検出など，
界面処理手法に共通する基盤である．

\textbf{近似精度}（付録~\ref{app:csf_delta_precision} に証明）：
Eikonal 条件 $|\bnabla\phi|=1$ のもとで，法線座標のテイラー展開と
$\delta_\varepsilon$ の偶関数性から：
\[
  \int_\Omega f\,|\bnabla\psi|\,dV - \int_\Gamma f\,dA_\Gamma
  = \Ord{\varepsilon^2\,\kappa^2\,\|f\|_{L^\infty}}
\]
$\varepsilon \sim \Delta x$ の設定では誤差は $\Ord{\Delta x^2}$ である．

\textbf{$\varepsilon$ が小さすぎる場合（$\varepsilon \ll \Delta x$）：}
界面が格子幅より薄くなり，密度・粘性の急激な変化が数値微分を困難にして
圧力 Poisson 方程式の係数行列が悪条件化し数値振動が生じる．

\textbf{$\varepsilon$ が大きすぎる場合（$\varepsilon \gg \Delta x$）：}
界面が人工的に厚くなり，曲率計算の精度が劣化するとともに
移流時の「にじみ」が積み重なって質量保存誤差が増大する．

\textbf{適正値（$\varepsilon \approx 1\text{--}2 \, \Delta x$）：}
界面が約 $2\text{--}4$ グリッドにわたって滑らかに変化することで，
物性の急変を緩和しつつ界面形状を解像度限界まで鮮明に保つ最適なバランスが実現される．

$C_\varepsilon$ を適切に設定することで，界面は中心遷移幅 $\approx 2\varepsilon$，
95\% 遷移帯 $\approx 6\varepsilon$ の滑らかな遷移層として維持される．

\medskip
\noindent\textbf{再初期化の離散化上の詳細．}
仮想時間刻み $\Delta\tau$ の安定性解析（双曲・放物型混合制約）と
設定指針（収束モニター $M(\tau)$，終了判定 $\delta_\tau$，$N_{\tau,\max}$）は
離散化レイヤに属する設計事項であり，§~\ref{sec:reinit} で定式化する．
ラベル \ref{warn:cls_dtau_stability} は §~\ref{sec:reinit} 側で定義される．
